\chapter{Operasyonel Guvenlik}

\section*{Giris}
Operasyonel guvenlik, guvenlik kontrollerinin gunluk hayatta olculebilir ve tekrarlanabilir sekilde isletilmesidir. NOC/SOC is birligi, olay mudahale ve tehdit avciligi bu bolumun omurgasini olusturur.

\section{Operasyon Merkezleri: NOC ve SOC Entegrasyonu}
\subsection{Ag Izleme ve Guvenlik Izleme Sinerjisi (ISOC)}
NOC ve SOC veri silolari kaldirilarak ortak telemetri ve ortak olay onceliklendirme modeli kurulmalidir.

\section{Yeni Nesil SOC Mimarisi ve Log Yonetimi}
\subsection{SIEM, SOAR Optimizasyonu ve Log Korelasyonu}
SIEM/SOAR optimizasyonunda use-case bakimi, log kalitesi, yanlis pozitif yonetimi ve playbook olgunlugu birlikte takip edilmelidir.

\section{Aktif Savunma ve Aldatma Teknolojileri}
\subsection{Bal Kupleri (Honeypot) ve Bal Aglari (Honeynet) Konumlandirmasi}
\subsection{Honeymonkey ve Istemci Tabanli Aldatma Sistemleri}
Aldatma altyapilari yalnizca tespit amacli degil, tehdit istihbarati ureten sensör katmani olarak konumlandirilmalidir.

\section{Siber Tehdit Istihbarati ve Tehdit Avciligi (Threat Hunting)}
\subsection{Ag Adli Bilisimi (Network Forensics) ile Anomali Avi}
Hipotez tabanli avcilikta ATT\&CK esleme, IOC/TTP zenginlestirme ve netflow/pcap analizi birlikte degerlendirilmelidir.

\section{Olay Mudahale (Incident Handling) ve Kriz Yonetimi}
\subsection{6 Adimli Olay Mudahale Yasam Dongusu (Hazirlik, Tespit, Sinirlandirma, Yok Etme, Kurtarma, Ders Cikarma)}
\subsection{Adli Bilisimde Delil Zinciri (Chain of Custody) Kurallari}
\subsection{Mavi Takim (Blue Team) Tatbikatlari ve Playbook Yonetimi}
Bu yasam dongu, teknik ekiplerin yani sira hukuk, iletisim ve is birimlerinin de dahil oldugu kurumsal kriz yonetimi modeliyle uygulanmalidir.
